\chapter{Proof Outline}

This blueprint documents the Lean package \texttt{LowRankUnivariateSOS}. It is
intended to be self-contained at the level of the formal proof strategy: the
optimization model, the algebraic reductions, the hgroup step, and the
factor-peeling transport are all stated here in the same order in which the
formal proof uses them.

The endpoint is the theorem \texttt{rankTwo\_no\_spurious\_socp} in
\texttt{RankTwoMain.lean}. The Lean statement is slightly stronger than the one
in the paper: it is proved for an arbitrary positive-definite bilinear form on
the polynomial space, not just for one fixed inner product.

This is the final formal claim.
\begin{theorem}\label{thm:main}
\uses{prop:abstract-socp,prop:setup}
\lean{LowRankUnivariateSOS.rankTwo_no_spurious_socp}
\leanok
For every sum-of-squares polynomial $p$ and every second-order critical point
$u=(u_1,u_2)$ of the quadratic-penalty objective,
\[
\sigma(u)-p = 0.
\]
Equivalently,
\[
\sigma(u)=p,
\]
so every SOCP is globally optimal.
\end{theorem}

\section{Optimization Model}

This definition fixes the objective and the exact FOCP/SOCP equations.
\begin{definition}\label{def:model}
\lean{LowRankUnivariateSOS.sigma2, LowRankUnivariateSOS.A, LowRankUnivariateSOS.residual, LowRankUnivariateSOS.objective, LowRankUnivariateSOS.IsSOS, LowRankUnivariateSOS.IsFOCP, LowRankUnivariateSOS.hessianTerm, LowRankUnivariateSOS.IsSOCP}
\leanok
The formal model is defined in \texttt{PolynomialModel.lean} and
\texttt{Socp.lean}. For a pair of polynomials $u=(u_1,u_2)$,
\[
\sigma(u)=u_1^2+u_2^2,\qquad
A_u(v)=u_1v_1+u_2v_2,\qquad
\operatorname{residual}(p,u)=\sigma(u)-p.
\]
The objective is
\[
f_p(u)=\langle \sigma(u)-p,\sigma(u)-p\rangle.
\]
The FOCP condition is
\[
\forall v,\qquad \langle A_u(v),\sigma(u)-p\rangle = 0.
\]
In Lean this is \texttt{IsFOCP}. The Hessian quadratic form is
\[
\operatorname{hessianTerm}(p,u,v)=
\langle \sigma(v),\sigma(u)-p\rangle + 2\|A_u(v)\|^2.
\]
The SOCP condition is the FOCP condition together with
\[
\forall v,\qquad 0 \le \operatorname{hessianTerm}(p,u,v).
\]
In the kernel directions used later in the proof,
\[
A_u(v)=0
\quad\Longrightarrow\quad
0 \le \langle \sigma(v),\sigma(u)-p\rangle.
\]
In Lean this full condition is \texttt{IsSOCP}.
\end{definition}

This is the abstract inequality package used at the end of the proof.
\begin{proposition}\label{prop:abstract-socp}
\uses{def:model}
\lean{LowRankUnivariateSOS.IsSOCP, LowRankUnivariateSOS.ImageOrthogonalResidual, LowRankUnivariateSOS.factored_h1_residual_eq_zero}
\leanok
Write
\[
r=\sigma(u)-p.
\]
Then the SOCP conditions give
\[
\forall q \in \operatorname{im}(A_{u_{\mathrm{img}}}),\qquad \langle q,r\rangle = 0
\]
for the transformed image used later in the proof, and
\[
\forall w \in \ker(A_u),\qquad 0 \le \langle \sigma(w),r\rangle.
\]
So if one can prove
\[
p = q_{\mathrm{img}} + \sum_i \sigma(w_i),
\qquad
q_{\mathrm{img}} \in \operatorname{im}(A_{u_{\mathrm{img}}}),
\qquad
w_i \in \ker(A_u),
\]
then
\[
0 \le \langle p,r\rangle
= \langle q_{\mathrm{img}},r\rangle + \sum_i \langle \sigma(w_i),r\rangle.
\]
Evaluating the FOCP condition at $v=u$ gives
\[
\langle \sigma(u),r\rangle = 0,
\]
hence
\[
f_p(u)=\langle r,r\rangle
= \langle \sigma(u)-p,r\rangle
= -\langle p,r\rangle \le 0.
\]
Since always $f_p(u)\ge 0$, it follows that
\[
f_p(u)=0,\qquad r=0.
\]
This is exactly the abstract content of
\texttt{factored\_h1\_residual\_eq\_zero}.
\end{proposition}

\section{Reduced Factorization}

This is the first algebraic reduction.
\begin{proposition}\label{prop:reduced-factorization}
\lean{LowRankUnivariateSOS.ReducedFactorization, LowRankUnivariateSOS.gcd_sigma_decomposition, LowRankUnivariateSOS.no_real_roots_sigma_reduced}
\leanok
Lean constructs a decomposition
\[
u = a \cdot u_0,
\qquad
\gcd(u_{0,1},u_{0,2})=1.
\]
In Lean this is the structure \texttt{ReducedFactorization} produced by
\texttt{gcd\_sigma\_decomposition}. The companion theorem
\texttt{no\_real\_roots\_sigma\_reduced} proves
\[
\forall x \in \mathbb R,\qquad \sigma(u_0)(x)\neq 0.
\]
So the common factor of $u_1,u_2$ is isolated in $a$, while the reduced pair
$u_0$ is coprime and its sum of squares has no real roots.
\end{proposition}

This separates the common factor into the good part and the peelable part.
\begin{proposition}\label{prop:factor-split}
\uses{prop:reduced-factorization}
\lean{LowRankUnivariateSOS.factor_peeling_decomposition}
\leanok
Assume $a\neq 0$. Then Lean produces a factorization
\[
a = g_{\mathrm{base}} \prod_i r_i
\]
such that
\[
\forall i,\qquad r_i \mid \sigma(u_0),
\qquad
\gcd(g_{\mathrm{base}},\sigma(u_0))=1.
\]
The peeled terminal factor is
\[
g=\operatorname{peelTerminalScale}(g_{\mathrm{base}},[r_1,\dots,r_n]),
\]
and satisfies
\[
\gcd(g,\sigma(u_0))=1.
\]
Informally, each bad factor $r_i$ is replaced by one copy of $X^2$, while the
coprime part $g_{\mathrm{base}}$ is kept unchanged.
\end{proposition}

\section{The hgroup Step and the \texorpdfstring{$h=1$}{h=1} Case}

This definition records the explicit image and kernel generators used in the
$h=1$ argument.
\begin{definition}\label{def:h1-machinery}
\lean{LowRankUnivariateSOS.scalePair, LowRankUnivariateSOS.perpPair, LowRankUnivariateSOS.inKerA_perpPair_scalePair, LowRankUnivariateSOS.inSigmaKerCone_sos_mul_sigma2, LowRankUnivariateSOS.inImageA_mul_of_isCoprime}
\leanok
The basic constructions are
\[
\operatorname{scalePair}(g,u)=(g u_1,g u_2),
\qquad
\operatorname{perpPair}(t,u)=(-t u_2,t u_1).
\]
They satisfy
\[
A_{g u}(\operatorname{perpPair}(t,u))=0,
\qquad
\sigma(\operatorname{perpPair}(t,u))=t^2 \sigma(u).
\]
Hence if
\[
s=q_1^2+q_2^2,
\]
then
\[
s\,\sigma(u)=\sigma(\operatorname{perpPair}(q_1,u))
            +\sigma(\operatorname{perpPair}(q_2,u)),
\]
so $s\,\sigma(u)$ lies in the cone generated by kernel squares. On the image
side, coprimality gives Bézout surjectivity:
\[
\gcd(u_1,u_2)=1
\quad\Longrightarrow\quad
\forall q\ \exists v,\ A_u(v)=q.
\]
Therefore
\[
\forall q\ \exists v,\ A_{g u}(v)=gq.
\]
\end{definition}

This is the algebraic division step behind the paper's hgroup argument.
\begin{lemma}\label{lem:hgroup}
\lean{LowRankUnivariateSOS.isSOS_mul, LowRankUnivariateSOS.hgroup_affine}
\leanok
If $p$ and $q$ are SOS polynomials and $\gcd(g,q)=1$, then there exist $s,t$
such that
\[
p = g\,t + s\,q,
\qquad
s \text{ is SOS}.
\]
Lean proves this by choosing Bézout coefficients
\[
ag+bq=1
\]
and defining
\[
s=b^2 p q,
\qquad
t=p(2a-a^2 g).
\]
Then
\[
g\,t+s\,q
=p\bigl(g(2a-a^2g)+b^2 q^2\bigr)
=p\bigl(ag(2-ag)+(bq)^2\bigr)
=p,
\]
using $bq=1-ag$. The SOS claim for $s$ is exactly
\[
b^2 \text{ is SOS},\quad p \text{ is SOS},\quad q \text{ is SOS}
\quad\Longrightarrow\quad
b^2 p q \text{ is SOS},
\]
formalized by \texttt{isSOS\_mul}.
\end{lemma}

This converts the hgroup identity into the paper's image-plus-cone condition.
\begin{proposition}\label{prop:h1}
\uses{def:h1-machinery,lem:hgroup}
\lean{LowRankUnivariateSOS.inImagePlusSigmaKerCone_of_h1_data}
\leanok
Assume
\[
\gcd(u_{0,1},u_{0,2})=1,
\qquad
p = g\,q + s\,\sigma(u_0),
\qquad
s \text{ is SOS}.
\]
Then
\[
gq \in \operatorname{im}(A_{g u_0}),
\qquad
s\,\sigma(u_0)\in
\operatorname{cone}\bigl(\sigma(\ker(A_{g u_0}))\bigr),
\]
and therefore
\[
p \in \operatorname{im}(A_{g u_0}) +
\operatorname{cone}\bigl(\sigma(\ker(A_{g u_0}))\bigr).
\]
This is the formal $h=1$ reduction used later by the abstract SOCP argument.
\end{proposition}

\section{Factor Peeling}

This definition packages the one-parameter family used in a single peeling
step.
\begin{definition}\label{def:peeling-witness}
\lean{LowRankUnivariateSOS.PeelingStepWitness, LowRankUnivariateSOS.RotatedPeelingData, LowRankUnivariateSOS.HasPeelingWitness, LowRankUnivariateSOS.PeelingPath, LowRankUnivariateSOS.PeelingSchedule}
\leanok
A peeling witness is a family $v_\eta$ satisfying
\[
A_u(v_\eta)=a_{\mathrm{const}},
\qquad
\sigma(v_\eta)=\eta^2 q_{\mathrm{quad}} + 2\eta q + \sigma_{\mathrm{const}},
\]
with
\[
q_{\mathrm{quad}} \in \operatorname{im}(A_{u_{\mathrm{img,prev}}}).
\]
So the quadratic coefficient is already controlled by the previous stage, while
the linear coefficient is the new target $q$. Lean stores the local step in
\texttt{PeelingStepWitness} and \texttt{RotatedPeelingData}. The recursive
transport is stored in
\[
\texttt{HasPeelingWitness},\qquad
\texttt{PeelingPath},\qquad
\texttt{PeelingSchedule}.
\]
\end{definition}

This is the one-factor orthogonality transport.
\begin{proposition}\label{prop:one-step-peeling}
\uses{def:peeling-witness,def:h1-machinery}
\lean{LowRankUnivariateSOS.rotatedPeelingData_of_factor_step, LowRankUnivariateSOS.dot_eq_zero_of_peeling_step, LowRankUnivariateSOS.image_orthogonal_of_peeling_step}
\leanok
Assume one quadratic factor step
\[
g_{\mathrm{next}}\,r = X^2 g_{\mathrm{start}},
\qquad
r \mid \sigma(u_0),
\qquad
r \mid g_{\mathrm{start}}.
\]
Let
\[
q\in \operatorname{im}(A_{g_{\mathrm{next}}u_0}),
\qquad
A_{g_{\mathrm{next}}u_0}(b)=q,
\]
and define
\[
\operatorname{quad}=\operatorname{perpPair}(1,g_{\mathrm{next}}u_0),
\qquad
\operatorname{const}=\operatorname{perpPair}(1,b),
\qquad
v_\eta=\eta\,\operatorname{quad}+\operatorname{const}.
\]
Then
\[
A_u(v_\eta)=A_u(\operatorname{const}),
\]
and
\[
\sigma(v_\eta)=
\eta^2 \sigma(\operatorname{quad})
 + 2\eta q
 + \sigma(\operatorname{const}).
\]
Moreover,
\[
\sigma(\operatorname{quad})\in \operatorname{im}(A_{g_{\mathrm{start}}u_0}),
\]
because $r$ divides both $\sigma(u_0)$ and $g_{\mathrm{start}}$. If the
previous image is already orthogonal to the residual $r_{\mathrm{res}}$, then
the SOCP inequality gives
\[
0 \le
\langle \sigma(v_\eta),r_{\mathrm{res}}\rangle
 + 2\langle A_u(v_\eta),A_u(v_\eta)\rangle
= 2\eta \langle q,r_{\mathrm{res}}\rangle + C
\]
for all real $\eta$. Hence
\[
\langle q,r_{\mathrm{res}}\rangle = 0.
\]
This is the content of \texttt{dot\_eq\_zero\_of\_peeling\_step}; the theorem
\texttt{image\_orthogonal\_of\_peeling\_step} upgrades it from one polynomial
$q$ to the entire next image.
\end{proposition}

This iterates the one-factor transport along the whole bad factor.
\begin{proposition}\label{prop:peeling}
\uses{prop:factor-split,prop:one-step-peeling}
\lean{LowRankUnivariateSOS.factor_peeling_schedule, LowRankUnivariateSOS.factor_peeling_path, LowRankUnivariateSOS.factor_peeling_image_orthogonal}
\leanok
Starting from
\[
a = g_{\mathrm{base}} \prod_i r_i,
\]
Lean builds a peeling schedule
\[
a u_0
\rightsquigarrow
(X^2 g_{\mathrm{base}} r_2 \cdots r_n)u_0
\rightsquigarrow \cdots \rightsquigarrow
g u_0.
\]
Orthogonality propagates at each arrow by Proposition~\ref{prop:one-step-peeling}.
Therefore the final theorem \texttt{factor\_peeling\_image\_orthogonal} gives
\[
\exists g\neq 0,\qquad
\gcd(g,\sigma(u_0))=1,
\qquad
\forall q\in \operatorname{im}(A_{g u_0}),\ \langle q,\sigma(u)-p\rangle = 0.
\]
This is the Lean form of the paper's factor-peeling conclusion.
\end{proposition}

\section{Final Assembly}

This packages exactly the algebraic data needed by Proposition~\ref{prop:abstract-socp}.
\begin{proposition}\label{prop:setup}
\uses{prop:reduced-factorization,prop:h1,prop:peeling}
\lean{LowRankUnivariateSOS.factor_peeling_certificate_step, LowRankUnivariateSOS.rankTwo_factored_h1_setup}
\leanok
Start from
\[
u=a u_0.
\]
If
\[
a=0,
\]
then the residual vanishes directly from the SOCP inequalities. If
\[
a\neq 0,
\]
then factor peeling gives
\[
\gcd(g,\sigma(u_0))=1,
\qquad
\forall q\in \operatorname{im}(A_{g u_0}),\ \langle q,\sigma(u)-p\rangle = 0.
\]
Applying Lemma~\ref{lem:hgroup} with $q=\sigma(u_0)$ gives
\[
p = g\,q' + s\,\sigma(u_0),
\qquad
s \text{ is SOS}.
\]
So \texttt{rankTwo\_factored\_h1\_setup} returns exactly the package
\[
u=a u_0,\qquad
p = g\,q' + s\,\sigma(u_0),\qquad
\gcd(g,\sigma(u_0))=1,
\qquad
\langle \operatorname{im}(A_{g u_0}),\sigma(u)-p\rangle = 0.
\]
\end{proposition}

\section{How the Proof Concludes}

This is the final substitution into the abstract SOCP inequality.

From Proposition~\ref{prop:setup} and Proposition~\ref{prop:h1}, the polynomial
$p$ has the form
\[
p = q_{\mathrm{img}} + \sum_i \sigma(w_i),
\qquad
q_{\mathrm{img}} \in \operatorname{im}(A_{g u_0}),
\qquad
w_i \in \ker(A_{a u_0}).
\]
So Proposition~\ref{prop:abstract-socp} applies and gives
\[
\sigma(u)-p=0.
\]
The theorem \texttt{rankTwo\_no\_spurious\_socp} is therefore the short chain
\[
\text{reduced factorization}
\Longrightarrow
\text{factor peeling}
\Longrightarrow
\text{hgroup} + h=1
\Longrightarrow
\text{image-plus-cone}
\Longrightarrow
\sigma(u)-p=0.
\]
